\section{Introduction}

Large language models can explain code errors in rich detail, describing why a variable is misnamed, why an algorithm handles edge cases incorrectly, or why a data structure choice leads to inefficiency. Yet these explanations are wrong more than half the time. Analysis of Self-Refine shows that 94\% of its failures trace to bad feedback: 33\% identify the wrong location and 61\% propose incorrect fixes \citep{madaan2023selfrefine}. Meanwhile, execution-based feedback from unit tests provides ground-truth correctness signals but offers no insight into \emph{why} code fails---only \emph{that} it fails.

This tension between \textbf{semantic richness} and \textbf{signal fidelity} has driven two parallel lines of research. Execution feedback methods like CodeRL \citep{le2022coderl} and RLTF \citep{liu2023rltf} use test pass/fail signals as reinforcement learning rewards, achieving strong results on code generation benchmarks. AI feedback methods like Self-Refine use LLM-generated critique for iterative refinement at test time. These approaches have been developed and evaluated in isolation, with no controlled comparison on the same model and benchmark setup.

We observe that fidelity and richness are \textbf{orthogonal dimensions}. High fidelity does not preclude high richness; the two properties simply have not been combined. This observation motivates a simple question: what if we use execution to \emph{verify} AI feedback rather than replace it?

We propose \textbf{Execution-Verified AI Feedback (EVAF)}, a mechanism that filters AI-generated code critiques through unit test execution before using them as training signals. The intuition is straightforward: AI suggestions that break tests are rejected; suggestions that maintain or improve correctness are accepted. The surviving feedback should retain AI's semantic explanations while inheriting execution's correctness guarantee.

EVAF operates as follows. Given a code generation problem where a baseline model produces incorrect output, an AI feedback model generates a critique and proposed fix. This fix is applied tentatively and executed against the problem's unit tests. If all tests pass, the critique is accepted as verified feedback; otherwise, it is rejected. The accept rate---the fraction of AI suggestions surviving verification---determines whether EVAF is viable: too low ($<$10\%) and EVAF degenerates to pure execution feedback; too high ($>$90\%) and the gating adds no value.

We implement EVAF using CodeT5-770M as the baseline generator and CodeLlama-7b-Instruct as the feedback model, targeting HumanEval's 164 problems. Our implementation comprises seven modules: data loading, baseline model wrapper, feedback model wrapper, execution gating with sandboxed test execution, metrics computation, visualization, and pipeline orchestration.

\textbf{Our contribution is primarily architectural and negative.} We report that the EVAF mechanism is implementable using existing tools, but our experimental validation failed: the pipeline crashed at 53\% completion (87/164 problems) before any metrics could be computed. The hypothesis remains untested rather than refuted---infrastructure failure, not theoretical refutation. We release our implementation and analysis to enable future work with proper infrastructure hardening.

The paper proceeds as follows. Section~\ref{sec:related} positions EVAF within execution and AI feedback literature. Section~\ref{sec:methodology} describes the EVAF mechanism and implementation. Section~\ref{sec:experiments} details our experimental setup. Section~\ref{sec:results} reports the incomplete results. Section~\ref{sec:discussion} discusses implications and limitations. Section~\ref{sec:conclusion} concludes with concrete next steps.
